\begin{topic}[id=pytorch_embedded,
             difficulty={Einstieg},
             type={Überblick + Einordnung},
             prerequisites={Grundlagen Python; Basisverständnis maschinelles Lernen},
             practice={optional},
             owner={Dozententeam},
             updated={2026-04-09},
             tags={ PyTorch, Datasets, Training, TorchScript, ExecuTorch, ONNX, Quantisierung }]{PyTorch – kompakte Einführung mit Embedded-Bezug}

\TopicDescription{PyTorch eignet sich hervorragend zum schnellen Prototyping. Du lernst Datasets/Dataloaders, `nn.Module`, Logging/Checkpointing und reproduzierbares Training. Dann werden Exportpfade (TorchScript/ONNX) und einfache Optimierungen betrachtet. Der Schwerpunkt liegt auf realistischen Abwägungen: Wie viel Genauigkeit kostet eine Optimierung und was bringt sie an Latenz/Größe? Mit diesem Rüstzeug kannst du Varianten fair vergleichen und geeignete Modelle für Edge-Szenarien vorbereiten.}

\TopicLeitfragen{\item Welche Minimal-Pipelines sind für kleine Projekte am effizientesten?
\item Wann lohnt sich TorchScript/ONNX/ExecuTorch?
\item Wie misst und dokumentiert man die Effekte einfacher Optimierungen?
}
\TopicScope{\item Kein tiefes Training großer Modelle.
\item Kein vollständiger Hardware-Portierungsleitfaden.
}
\TopicPractice{Zwei Modellvarianten exportieren (TorchScript/ONNX) und Inferenzzeiten auf CPU vergleichen.}

\TopicKeywords{PyTorch, Datasets, Training, TorchScript, ExecuTorch, ONNX, Quantisierung}

\nocite{pytorchTutorials,execuTorchDocs}
\end{topic}
